% ID: SEL-D0060-OA-E04
% ESTADO: cotejada. Segunda pasada visual y borrador compilado revisados.
% Fuente: Selectividad/sel_2011_matematicas_aplicadas/223A-EXAMEN 3.pdf
% SHA-256: 5998699e0b2992b988a8e0ed0d2e3581dd8fac52d6f9802587ebd40fee5ef608
% Página: 1; opción A; ejercicio original 4.

El peso de los adultos de una determinada población sigue una distribución Normal de media 70 kg y desviación típica 16 kg. Si elegimos, al azar, muestras de tamaño 4,
\begin{enumerate}[label=\alph*)]
\item \textbf{(0.5 puntos)} \textquestiondown{}cuál es la distribución de la media muestral?
\item \textbf{(1 punto)} \textquestiondown{}cuál es la probabilidad de que el peso medio de una de esas muestras esté comprendido entre 65 y 72 kg?
\item \textbf{(1 punto)} \textquestiondown{}cuál es la probabilidad de que ese peso medio sea menor que 70kg?
\end{enumerate}
